% Chapter Template

\chapter{Ensayos y resultados}

\label{Chapter4}

En este capítulo se presentó el plan de evaluación del sistema, las pruebas automatizadas que respaldaron la reproducibilidad del pipeline implementado en el capítulo~\ref{Chapter3} y los resultados obtenidos al ejecutar el prototipo sobre los videos disponibles para el trabajo.

%----------------------------------------------------------------------------------------
%	SECTION 1
%----------------------------------------------------------------------------------------

\section{Pruebas automatizadas del pipeline}
\label{sec:pruebas-automatizadas}

Durante el desarrollo se mantuvo una batería de pruebas unitarias sobre componentes aislados del flujo analítico: calibración ChArUco, detección de fase estática, impulsos de audio, métricas espectrales y de tapping, ángulos derivados y escritores de exportación. Los módulos bajo \texttt{tests/unit/} validaron comportamientos con señales sintéticas o entradas controladas, de modo que las definiciones matemáticas descritas en el capítulo~\ref{Chapter3} pudieran contrastarse con salidas esperadas. Se incluyó además una prueba de integración de humo (\texttt{tests/integration/test\_cli\_smoke.py}) que ejecutó el comando de sesión sobre un video de muestra del repositorio y verificó la generación de \texttt{tracking.csv}, \texttt{metrics.csv} y \texttt{summary.json} con esquemas estables.

Esas pruebas no sustituyeron la validación clínica formal, excluida del alcance del trabajo, pero demostraron que el ensamblado descrito en el capítulo de diseño era reproducible y detectaba regresiones en el flujo de extremo a extremo.

%----------------------------------------------------------------------------------------
%	SECTION 2 — plan de evaluación, resultados y análisis (por completar)
%----------------------------------------------------------------------------------------
